% Summary of cole1993: Feature percolation

\begin{abstract}
\begin{enumerate}
  \item Introduces the core Feature Percolation Principles and one Feature Percolation Condition: daughters percolate their features to their mothers, the head daughter wins ties, and complements of lexical heads don't percolate their features
  \item Argues that the FPP explains Long-distance Reflexives and blocking effects in Chinese, Korean, Icelandic, and Italian; Non-C-Commanding Antecedents of Reflexives in Chinese and Korean, and non-c-commanding Honorific Agreement in Korean.
\end{enumerate}
\end{abstract}

\section*{Feature Percolation Principles}

\begin{enumerate}
  \item The features of the mother node and the features of the daughter nodes will be identical.
  \item If the features of the two daughter nodes conflict, the mother node will have the features of the head node.
\end{enumerate}

\section*{Feature Percolation Condition}

\begin{enumerate}[start=3]
  \item No feature can percolate out of a lexical complement structure.
\end{enumerate}

\section*{Long-distance reflexives}

\begin{enumerate}
  \item Zhangsan\textsubscript{i} thinks Lisi knows Wangwu likes self\textsubscript{i}.
\end{enumerate}

\subsection*{Blocking effects}

In Chinese and Korean, if an intermediate subject differs in phi-features (not sure exactly what this means) from a lower subject, long-distance reflexives are blocked.

This isn't true in Italian or Icelandic:

\begin{enumerate}[start=2]
  \item Gianni\textsubscript{i} suppone che tu sia innamorato della propria\textsubscript{i} moglie.
  \item Anna\textsubscript{i} believes you-ACC have-INF betrayed self\textsubscript{i}.
\end{enumerate}

\section*{Non-c-commanding antecedents of reflexives}

\begin{enumerate}[start=4]
  \item Zhangsan\textsubscript{i} 's pride hurt self\textsubscript{i}.
  \item John\textsubscript{i} 's father self\textsubscript{i} ruined.
\end{enumerate}

Antecedents for reflexives must also be animate in Chinese and Korean.

Plus, an animate head percolates instead of the possessor (in Chinese but not Korean, see (5) - maybe Korean `father' isn't animate?):

\begin{enumerate}[start=6]
  \item Zhangsan\textsubscript{i} 's father\textsubscript{j} to self\textsubscript{*i/j} no confidence
\end{enumerate}

'Zhangsan's father has no confidence in himself'

\section*{Honorific agreement in Korean}

In Korean only the subject triggers an HON affix on the verb.

If the subject is not honorific, but has an honorific possessor, this still happens.

Sometimes (it's less felicitous but still possible), if the honorific is in a relative clause describing the subject, this can still work:

\begin{enumerate}[start=7]
  \item The clothes which the teacher is wearing are very nice-HON.
\end{enumerate}

\section*{Percolation from complements (banned)}

\begin{enumerate}[start=8]
  \item *[I\textsubscript{i} scold him\textsubscript{j}] to self\textsubscript{i/*j} not have advantage
\end{enumerate}
`that I scolded him did myself no good'

The feature of `him' should percolate instead of the feature of `I', because it should win out as the feature of the head when the two merge, but it doesn't.

They argue that complements of lexical categories (like V in this case) never project.

They get around the fact that pied-piping happens with objects of prepositions by referring to Halbert 1990, who argues that `objects' of prepositions are actually specifiers not complements.
